\chapter{变量与类型}
\label{ch:03}

程序中的数据需要有名字，也需要有解释它的类型。同一个数字可能表示分数、计数或容量，但合法范围、运算方式和是否允许修改未必相同。选类型先从数据含义出发，再看写法是否简洁。

本章依次讨论声明与更新、字面量定型、重载观察和显式转换。每个例程都追踪“声明时是什么类型、执行后是什么值”这两个问题，避免把类型推导理解成运行时随意换类型。

\section{声明、赋值与推导}

从成绩变量开始。声明给它一个名字和 int64 类型，初始化把 70 作为第一份值放进去；后面的赋值改变当前值，而不重新选择变量类型。把这两个动作分开，就能理解同一个 score 为什么先是 70，后来是 75。

如果右侧已经给出足够信息，可以用 let 让编译器推导类型。需要表达“不再通过这个绑定修改”时再使用 const。推导和只读解决不同问题；普通 const 也不自动满足所有编译期常量求值要求，例如数组维度还要符合相应规则。

\par
\medskip
\begin{lstlisting}[style=chali]
void main()
{
    int64 score = 70;
    score += 5;
    let count = 2;
    const int64 limit = 100;
    println("{} {} {}", score, count, limit);
}
\end{lstlisting}
\par
\medskip

\par
\medskip
\begin{lstlisting}[style=out]
75 2 100
\end{lstlisting}
\par
\medskip

\subsection*{代码详解}

\texttt{int64 score =\allowbreak{} 70} 同时给出类型、变量名与初值。下一行的 \texttt{+=\allowbreak{} 5} 读取当前 score，加上 5 后写回，所以打印前它已变为 75。\texttt{let count =\allowbreak{} 2} 省略显式类型，让编译器从初始化表达式推导；省略类型不等于省略初始化。

\texttt{const int64 limit =\allowbreak{} 100} 表明这个绑定在后续不能被普通赋值修改。输出格式串中三个 \texttt{\{\}} 分别接收 score、count、limit，按参数位置得到 75、2、100；花括号本身不是输出数据。先分清声明、更新和只读约束，再阅读最终输出，会比只记住 let 的拼写更有用。

\section{数值、布尔与文本}

整数类型按位宽和有无符号区分，float32／float64 表示浮点数。bool 用于条件，char 表示字符值，String 保存 UTF-8 字节序列。不要把字符个数、字节数与整数取值范围当成同一概念。浮点计算存在表示误差；精确十进制、大整数等需求查相应标准库。

无其他约束的整数字面量通常按 int64 处理；后缀、目标类型与重载选择会影响具体类型。下例用重载观察浮点字面量，并展示不同整数范围。

\par
\medskip
\begin{lstlisting}[style=chali]
int64 floatingKind(float32 value)
{
    return 32;
}

int64 floatingKind(float64 value)
{
    return 64;
}

uint64 unsignedValue(uint64 value)
{
    return value;
}

void main()
{
    let one = 1u;
    let single = 1.0f;
    uint8 narrow = 255u;
    uint64 maximum = 18446744073709551615u;
    println("{} {} {} {}", unsignedValue(one), floatingKind(single),
            floatingKind(1.0), narrow);
    println("{}", maximum);
}
\end{lstlisting}
\par
\medskip

\par
\medskip
\begin{lstlisting}[style=out]
1 32 64 255
18446744073709551615
\end{lstlisting}
\par
\medskip

\subsection*{代码详解}

两个 floatingKind 同名但形参类型不同，构成重载；它们直接返回 32 或 64，用来观察编译器选择了哪一个函数，而不是计算浮点数位数。\texttt{1.\allowbreak{}0f} 带 float32 后缀，因此 single 选择前者；没有 f 的 \texttt{1.\allowbreak{}0} 选择 float64 版本。

\texttt{1u} 与后面的无符号常量展示无符号定型。\texttt{uint8 narrow =\allowbreak{} 255u} 还有目标类型提供的范围约束，不能因为带 u 就放入任意大的整数。maximum 使用 uint64 能容纳的最大值，第二次打印核对整个十进制值没有被当作有符号整数。第一行四项依次来自无符号恒等函数、两个重载结果和 narrow，而不是四次不同的转换。

\section{转换写在表示边界}

将 int64 保存为 int32 可能丢失范围，因此缩窄转换需要显式表达。可使用 \texttt{(\allowbreak{}int32)\allowbreak{}value} 或 \texttt{cast<\allowbreak{}int32>\allowbreak{}(\allowbreak{}value)\allowbreak{}}。允许的隐式转换和字面量范围检查见语言参考；显式强转也不是业务范围校验，输入仍应按应用规则验证。

\par
\medskip
\begin{lstlisting}[style=chali]
void main()
{
    int64 original = 42;
    int32 narrowed = cast<int32>(original);
    println("{}", narrowed);
}
\end{lstlisting}
\par
\medskip

\par
\medskip
\begin{lstlisting}[style=out]
42
\end{lstlisting}
\par
\medskip

\subsection*{代码详解}

original 声明为 int64，narrowed 声明为 int32。\texttt{cast<\allowbreak{}int32>\allowbreak{}(\allowbreak{}original)\allowbreak{}} 把缩窄意图写在赋值右侧，结果再交给 narrowed；它没有把 original 的声明类型改成 int32。

42 落在目标类型范围内，因此本例打印 42。这个输入只能展示显式转换的写法，不能证明任意 int64 都适合转换。处理外部数据时，应先按业务与目标类型范围检查，再决定是否转换，不能把 cast 当成输入验证函数。

\section{作用域与后续使用}

同一个名字在不同作用域可以表示不同绑定。读取变量前需要完成初始化；移动后的可用性另由所有权规则约束。下一步将这些值放入函数参数和返回类型中，观察类型如何约束调用。

\subsection*{小结}

读变量时按名字、类型、当前值、可写性四项核对。let 只省略可推导的类型书写，const 限制修改，cast 表达转换意图；三者解决不同问题。类型合法后，业务还可能要求年龄或成绩落在更小范围内。

下一章把这些约定放到函数边界。形参类型说明接受什么，返回类型说明交回什么，调用过程仍遵守本章的兼容与转换规则。

\section{习题}

\begin{enumerate}

\item 把成绩改成 float64 并输出。

\item 比较 let 和 Any。

\item int64 到 int32 为什么要考虑范围？

\item const 是否等同于编译期执行？

\end{enumerate}

\subsection*{参考解答与验收要点}

\begin{enumerate}

\item 改变类型和初始化式，核对小数输出。

\item let 推导具体类型；Any 是显式的类型擦除值。

\item 目标表示范围较小；先按业务要求检查，再决定转换。

\item 不同。常量求值另有允许的表达式和计算预算。

\end{enumerate}

下一章：\hyperref[ch:04]{第 4 章}。

\section{番外：小杯装不下鲲鹏}

改写自《庄子·逍遥游》。

北冥有数，其名为鲲。鲲之大，不知其几千亿也；化而为数组，其名为鹏。鹏之大，不知其几千万项也。怒而分配，其请求若垂天之云。

《齐谐》者，志怪者也。《谐》之言曰：“鹏之申请内存也，水击三千行，抟报错而上者九万层，去以六个月之预算者也。”

蜩与学鸠笑之曰：“我决起而算，抢一个 \texttt{int8} 而止，时则不至，而控于地而已矣，奚以九万亿为？”它们用一只八位小杯接住账房递来的总数，杯沿顿时溢得满地都是。学鸠擦干桌面，宣布问题已解决，因为屏幕上又出现了一个小数字。

朝菌不知晦朔，蟪蛄不知春秋，窄整数不知业务总额。小知不及大知，小类型装不下大输入，奚以知其然也？看一下范围便知。

鹏说：“且夫水之积也不厚，则其负大舟也无力。你连数值上界都没写，怎么负得动年终报表？”

蜩取出一张显式转换，说自己已经签字，溢出应当自理。庄子叹道：“此亦类型之辩也。签字说明你想倒进去，不说明杯子忽然长成了海。”

\subsection*{技术对照}

数值类型同时规定表示范围和运算行为。显式转换表明转换意图，并不自动证明输入落在目标范围内；先检查边界，再选择表示类型。
